\chapter{Strategi dan Portofolio API}

\section*{Tujuan Pembelajaran}
\begin{itemize}
  \item Menurunkan kandidat API dari tujuan strategis dan kapabilitas bisnis organisasi.
  \item Memprioritaskan portofolio API dengan kriteria nilai, usaha, dan risiko.
  \item Menyusun kanvas produk API untuk API prioritas.
\end{itemize}

\section{Pendahuluan}

% TODO
% Cakupan: setelah kondisi saat ini terpetakan, Medikara harus memutuskan API mana yang dibangun
% lebih dulu dengan anggaran terbatas, dan siapa yang bertanggung jawab atasnya.

\section{Dari Strategi Platform ke Strategi API}

% TODO
% Cakupan: hubungan keputusan strategi platform (C1) dengan pilihan API; tujuan strategi API;
% prinsip API-first sebagai kebijakan organisasi; kesalahan umum berupa membuka API tanpa tujuan.
% Kasus Medikara: turunan strategi "Satu Pasien, Satu Rekam Medis" menjadi tujuan API.

\section{Menurunkan Kandidat API dari Kapabilitas Bisnis}

% TODO
% Cakupan: memakai capability map dan value stream dari IT40513 sebagai masukan; kapabilitas yang
% perlu dibuka ke mitra; perbedaan daftar API berbasis kapabilitas dan berbasis sistem.
% Kasus Medikara: kapabilitas penjadwalan, rekam medis, hasil laboratorium, dan penagihan.

\section{Batas Domain dan Lapisan API}

% TODO
% Cakupan: batas API mengikuti batas domain bisnis; lapisan API sistem, proses, dan pengalaman
% beserta trade-off; Domain-Driven Design hanya disinggung karena dibahas rinci di IT40713.
% Pendalaman teknis: memetakan kandidat API ke batas domain dan menandai API yang melintasi lebih
% dari satu domain.

\section{Lanskap API}

% TODO
% Cakupan: lanskap API sebagai kumpulan API milik organisasi; delapan aspek lanskap menurut Medjaoui
% dkk. (variety, vocabulary, volume, velocity, vulnerability, visibility, versioning, volatility);
% katalog API sebagai alat melihat lanskap.
% Kasus Medikara: lanskap API yang tumbuh liar dari lima rumah sakit hasil akuisisi.

\section{Kanvas Produk API}

% TODO
% Cakupan: elemen kanvas (pengguna, masalah yang diselesaikan, nilai, data, kapabilitas, ukuran
% keberhasilan, risiko, pemilik); cara mengisi kanvas bersama tim campuran.
% Kasus Medikara: kanvas API janji temu sebagai contoh lengkap.

\section{Prioritisasi Portofolio}

% TODO
% Cakupan: kriteria nilai, usaha, dan risiko; matriks nilai dan usaha; keputusan membangun sendiri,
% membeli, atau memakai API pihak lain; peran pemilik produk API; business case menyeluruh dibahas
% rinci di IT40513.
% Kasus Medikara: memilih tiga API dari lima kandidat dengan anggaran terbatas.

\section{Aktivitas: Portofolio API Target}

% TODO
% Tujuan aktivitas: menetapkan portofolio API target yang berprioritas dan memiliki pemilik.
% Peran dalam tim: anggota berlatar bisnis menilai nilai bagi pengguna dan mitra; anggota berlatar
% teknis menilai usaha, ketergantungan sistem, dan risiko.
% Langkah: turunkan kandidat dari tujuan strategis dan kapabilitas; isi matriks nilai dan usaha;
% pilih dua atau tiga API teratas; isi kanvas produk API untuk masing-masing; tentukan pemiliknya.
% Luaran dan tonggak capstone: portofolio API target dan kanvas produk API, masuk ke C3.
% Berkas sources yang direncanakan: sources/chapter08/kanvas-produk-api.md,
% sources/chapter08/matriks-prioritas-api.csv, sources/chapter08/contoh-kanvas-janji-temu.md

\section{Ringkasan}

% TODO

\section{Latihan}

\begin{enumerate}
  \item Jelaskan mengapa daftar API yang diturunkan dari daftar sistem berbeda dengan daftar API yang diturunkan dari kapabilitas bisnis.
  \item Misalkan Medikara hanya punya anggaran untuk tiga API tahun depan. Susun kriteria prioritisasi, lalu terapkan pada lima kandidat API.
  \item Tentukan siapa pemilik produk yang tepat untuk API janji temu Medikara, lalu jelaskan tanggung jawabnya.
\end{enumerate}
